\chapter{Datengrundlage}
\label{kap:daten}

Alle Parameter der folgenden Kapitel berufen sich auf Messungen an dem
verfügbaren Quellmaterial. Dieses Kapitel dokumentiert diese Messungen. Sie sind
über das Skript \code{src/analyze\_data.py} reproduzierbar; die Rohwerte liegen in
\code{docs/datensichtung.json}.

\section{Annotationsdefinition}
\label{sec:annotation}

Vor jeder Messung steht die Frage, was überhaupt ein Objekt ist. Aus der Sichtung
der beiden handannotierten Realbilder ergibt sich eine Definition, die die
gesamte Pipeline prägt:

\begin{quote}
Annotiert ist ausschließlich die \textbf{sichtbare Stirnfläche} eines
Rundrohrprofils -- der Kreisring, den man am Ende eines Rahmenrohres sieht:
heller verzinkter Kranz, dunkle Bohrung. Rohrschäfte in Seitenansicht,
Gitterknoten und dunkle Löcher im Maschendraht sind \textbf{keine} Objekte,
obwohl sie ähnlich aussehen. Die Box umschließt den \textbf{sichtbaren} Teil
(modale Box), nicht den geschätzten vollständigen Kreis.
\end{quote}

Beide Festlegungen sind für die Synthese folgenreich. Die erste bestimmt, welche
Strukturen als harte Negative benötigt werden (Abschnitt~\ref{sec:negative}); die
zweite bestimmt, wie die Boxen im Generator berechnet werden. Die Unterscheidung
zwischen modaler und amodaler Box ist in der Literatur etabliert
\parencite{li2016amodal,zhu2017amodal}: Amodale Boxen umschließen die geschätzte
vollständige Objektausdehnung einschließlich verdeckter Teile. Würde der
Generator amodale Boxen erzeugen, während die Realbilder modal annotiert sind,
entstünde ein systematischer Größenversatz zwischen Trainings- und
Bewertungslabels -- ein Verstoß gegen Anforderung~A5 aus
Abschnitt~\ref{sec:anforderungen}.

\section{Objektausschnitte (Cutouts)}
\label{sec:cutouts}

Verfügbar sind neun aus Realaufnahmen ausgeschnittene Rohrprofile
(Tabelle~\ref{tab:cutouts}).

\begin{table}[htbp]
\centering
\caption[Die neun Objektausschnitte]{Die neun verfügbaren Objektausschnitte.
Die Ovalität ist das Verhältnis der Halbachsen des angepassten Ellipsenmodells
(Abschnitt~\ref{sec:matting}).}
\label{tab:cutouts}
\small
\begin{tabularx}{\textwidth}{@{}lrlcrX@{}}
\toprule
\textbf{Datei} & \textbf{Größe} & \textbf{Modus} & \textbf{Bohrung} &
\textbf{Ovalität} & \textbf{Bewertung} \\
\midrule
\code{cutted1.png}  & $38\times32$ & RGBA & dunkel & 0{,}86 & brauchbar; Nachbarrohr unten rechts im Ausschnitt \\
\code{cutted\_4.png} & $30\times30$ & RGBA & hell   & 1{,}14 & brauchbar; deutlich oval, Kranz einseitig hell \\
\code{cutted\_5.png} & $43\times38$ & RGBA & dunkel & 0{,}94 & bestes Exemplar, größte Auflösung \\
\code{cutted\_6.png} & $29\times30$ & RGBA & hell   & 1{,}14 & brauchbar; kontrastarm \\
\code{cutted\_7.png} & $41\times35$ & RGBA & dunkel & 0{,}87 & sehr sauber, kräftiger Kranz \\
\code{cutted\_8.png} & $39\times35$ & RGBA & dunkel & 0{,}83 & schwächstes Exemplar: Kranz kontrastarm \\
\code{cutted\_9.png} & $28\times30$ & RGBA & hell   & 0{,}97 & brauchbar \\
\code{cuuted2.png}  & $36\times33$ & RGBA & dunkel & 0{,}91 & brauchbar, bräunlicher Farbstich \\
\code{cuuted3.png}  & $36\times34$ & RGBA & dunkel & 0{,}92 & brauchbar \\
\bottomrule
\end{tabularx}
\end{table}

\absatz{Befund 1 -- der Alphakanal ist unbrauchbar.} Alle neun Dateien sind zwar
im RGBA-Format, der Alphakanal ist jedoch durchgängig auf 255 gesetzt (gemessen:
$\min(\alpha) = \max(\alpha) = 255$ für alle neun). Es handelt sich also um
rechteckige Bildausschnitte, nicht um freigestellte Objekte. Der Anteil der
Fläche, der gar nicht zum Objekt gehört, liegt bei 13--35\,\%. Direktes
Einkopieren erzeugte damit genau die Signatur~S1 aus
Abschnitt~\ref{sec:shortcut}. Das Freistellen ist deshalb keine Vorbereitung,
sondern eine eigene Pipelinestufe (Abschnitt~\ref{sec:matting}).

\absatz{Befund 2 -- zwei Erscheinungsformen.} Sechs Ausschnitte zeigen eine
\emph{dunkle} Bohrung (offenes Rohr, Blick ins Innere), drei eine \emph{helle}
(Sonnenlicht fällt in die Röhre, oder das Ende ist gestaucht bzw.\ verschlossen).
Beide Formen kommen in den Realbildern vor. Die Pipeline zieht gleichverteilt aus
allen neun; eine Gewichtung wäre ohne Häufigkeitszählung in den Realbildern nicht
belegbar und würde damit Anforderung~A6 verletzen.

\absatz{Befund 3 -- Ovalität ist im Quellmaterial vorhanden.} Die aus dem
Ellipsenfit gewonnenen Achsverhältnisse streuen zwischen 0,83 und 1,14. Ursachen
sind Schrägansicht (Projektion des Kreises auf eine Ellipse) und tatsächlich
gequetschte Rohre. Dieser Befund rechtfertigt die Ovalisierungstransformation
der Pipeline (Abschnitt~\ref{sec:instanzgeometrie}) -- allerdings nur der
Größenordnung nach; der tatsächlich verwendete Anteil wird aus der
Seitenverhältnisstatistik der Realbilder abgeleitet.

\absatz{Bewertung.} Ausreichend, aber knapp. Neun Instanzen sind sehr wenig; im
erzeugten Trainingssatz erscheint jede rund 15\,100-mal als gelabeltes Objekt
(136\,384 Labels / 9). Das Risiko, dass das Netz die neun konkreten Texturmuster
memoriert, ist real. Gegenmaßnahmen in der Pipeline sind starke photometrische
Streuung pro Instanz, Rotation über den vollen Kreis, Ovalisierung, Dellen und
Perspektive -- sowie die bewusste Wahl eines kleinen Modells
(Abschnitt~\ref{sec:trainingsaufbau}).

\section{Hintergründe}
\label{sec:hintergruende}

\begin{table}[htbp]
\centering
\caption[Die sieben Hintergrundaufnahmen]{Die sieben Hintergrundaufnahmen.
\enquote{Clean Plate} bezeichnet eine Aufnahme der Zieldomäne \emph{ohne}
sichtbare Stirnflächen. Rohr-\Oe{} ist der gemessene Rohrdurchmesser in Pixeln;
er verankert den Maßstab (Befund~6). Das Gewicht ist das Ziehgewicht im
Generator.}
\label{tab:hintergruende}
\small
\begin{tabularx}{\textwidth}{@{}lrlrrX@{}}
\toprule
\textbf{Datei} & \textbf{Größe} & \textbf{Typ} & \textbf{Rohr-\Oe} &
\textbf{Gewicht} & \textbf{Inhalt} \\
\midrule
\code{test\_cutted.png}  & $1513\times959$  & Clean Plate & 16 & 6 & Stapelwand, exakt die Zieldomäne \\
\code{test2\_cutted.jpg} & $1706\times1108$ & Clean Plate & 18 & 6 & Stapelwand, Gittermatten, Laub \\
\code{test4\_cutted.png} & $1517\times2788$ & Clean Plate & 21 & 6 & größte Fläche, hochkant, Laub am Boden \\
\code{test5\_cutted.jpg} & $945\times552$   & Clean Plate & 20 & 6 & Aufsicht, Absperrband, Gras \\
\code{bg\_3.jpg} & $1024\times576$ & Fremdszene & -- & 1 & Kieswerk, Schotter, Himmel \\
\code{bg\_4.jpg} & $800\times600$  & Fremdszene & -- & 1 & Betonplatten, Sandsäcke, Personen \\
\code{bg\_5.jpg} & $800\times600$  & Fremdszene & -- & 1 & Betondeck mit Maschendrahtzaun \\
\bottomrule
\end{tabularx}
\end{table}

\absatz{Befund 4 -- vier echte Clean Plates.} Die vier \code{test*}-Dateien sind
Aufnahmen derselben Objektart, aus denen die Stirnflächen herausgeschnitten
wurden: Wände aus Rohrschäften \emph{ohne} Stirnflächen. Sie sind der wertvollste
Teil des Materials, weil sie drei Dinge gleichzeitig liefern: den photometrisch
exakt passenden Kontext für das Einkopieren, das Rohmaterial für die
Okkluder-Bibliothek (Abschnitt~\ref{sec:okkluderbib}) und echte
Negativbeispiele. Sie erhalten deshalb sechsfaches Ziehgewicht. In der
Terminologie von Gleichung~\eqref{eq:bendavid} senken sie unmittelbar den
Divergenzterm $d_{\mathcal{H}\Delta\mathcal{H}}$: Der Hintergrund ist nicht
\enquote{ähnlich} der Zieldomäne, er \emph{ist} sie.

\absatz{Befund 5 -- die drei Fremdszenen sind hartes Negativmaterial.}
\code{bg\_5.jpg} enthält einen Maschendrahtzaun mit rautenförmigen Öffnungen --
genau die Struktur, die nach der Annotationsdefinition kein Objekt ist und
dennoch als Falschpositiv droht. \code{bg\_3.jpg} liefert Schotter, also viele
dunkle, runde Kiesel. Diese Bilder sind primär als Negativbilder wertvoll, nicht
als Objektträger (Abschnitt~\ref{sec:negative}).

\absatz{Befund 6 -- der Rohrdurchmesser verankert den Maßstab.} Aus den Clean
Plates lässt sich der Rohrdurchmesser messen (16/18/21/20\px{}). Eine Stirnfläche
hat zwangsläufig \emph{denselben} Durchmesser wie der Schaft daneben -- es ist
dasselbe Rohr. Der Generator zieht deshalb nicht Zoomfaktor und Objektgröße
unabhängig, sondern koppelt sie (Abschnitt~\ref{sec:hintergrundstufe}). Dies ist
das erste Beispiel für eine Parameterentscheidung, die nicht aus einer Statistik,
sondern aus einer physikalischen Randbedingung folgt (Anforderung~A3).

\absatz{Bewertung.} Gut, mit einer klaren Lücke: Alle vier Clean Plates zeigen
dieselbe Stapelart bei bedecktem bis diffusem Licht -- kein harter Sonnenschein,
keine Nässe, kein Schnee, kein Nachtbild. Diese Lücke kann die Pipeline nur
teilweise über photometrische Augmentierung schließen; sie bleibt die größte
verbleibende Domänenlücke (Abschnitt~\ref{sec:grenzen}).

\section{Validierungs- und Testdaten}
\label{sec:valtest}

Zur Bewertung stehen zwei vollständig handannotierte Realaufnahmen zur
Verfügung. Sie werden strikt getrennt verwendet: \code{val} dient der
Modellauswahl (Auswahl von \code{best.pt}, Early Stopping),
\code{test} ausschließlich der Endbewertung und wird von Training und
Modellauswahl nie berührt.

\begin{table}[htbp]
\centering
\caption[Die realen Bewertungsdaten]{Die realen Bewertungsdaten.}
\label{tab:valtest}
\small
\begin{tabular}{@{}llrr@{}}
\toprule
\textbf{Split} & \textbf{Bild} & \textbf{Auflösung} & \textbf{Objekte} \\
\midrule
\code{val}  & \code{PXL\_20251204\_073853552.jpg} & $2459\times1174$ & 88 \\
\code{test} & \code{test5.jpg}                    & $1600\times900$  & 86 \\
\bottomrule
\end{tabular}
\end{table}

\subsection{Objektstatistik}

\begin{table}[htbp]
\centering
\caption[Gemessene Objektstatistik]{Gemessene Objektstatistik der realen
Labels. Die Kenngröße $\sqrt{\text{Fläche}}$ dient im Folgenden als
Objektdurchmesser.}
\label{tab:objektstatistik}
\small
\begin{tabular}{@{}lll@{}}
\toprule
\textbf{Größe} & \textbf{\code{val}} & \textbf{\code{test}} \\
\midrule
Breite (px) & 19{,}9--38{,}5, Median 29{,}3 & 9{,}3--33{,}3, Median 23{,}8 \\
Höhe (px) & 15{,}7--34{,}2, Median 27{,}0 & 6{,}4--30{,}5, Median 19{,}7 \\
$\sqrt{\text{Fläche}}$ (px) & 20{,}6--35{,}8, \textbf{Median 27{,}4} & 12{,}2--31{,}9, \textbf{Median 19{,}6} \\
Seitenverhältnis $w/h$ & p10 0{,}97, Median 1{,}08, p90 1{,}31 & p10 0{,}85, Median 1{,}09, p90 1{,}22 \\
Objektfläche gesamt & 2{,}48\,\% des Bildes & 2{,}69\,\% \\
am Bildrand angeschnitten & 0\,\% & 5{,}8\,\% \\
\bottomrule
\end{tabular}
\end{table}

\absatz{Befund 7 -- sehr kleine Objekte.} Bei 20--36\px{} Kantenlänge in Bildern
von 1600--2459\px{} Breite handelt es sich um Kleinobjekte im Sinne der
COCO-Konvention \parencite{lin2014coco}. Nach Gleichung~\eqref{eq:deff} schrumpfen
sie bei \code{imgsz=640} auf 5--9\px{} und liegen damit unter der Auflösung des
feinsten Detektionskopfes (Stride 8). Die Konsequenz für die Bewertung ist
zwingend: Inferiert wird bei nativem Maßstab, und die Inferenzauflösung wird
systematisch untersucht (Abschnitt~\ref{sec:imgsz}).

\absatz{Befund 7b -- die Seitenverhältnisse weichen deutlich vom Quadrat ab.}
Über beide Splits gerechnet liegt das p10--p90-Intervall des Seitenverhältnisses
bei 0,85--1,31. Über die Hälfte der echten Boxen weicht also um mindestens
5\,\% vom Quadrat ab. Dieser Wert ist die Belegbasis für den Anteil der
Ovalisierung (Abschnitt~\ref{sec:instanzgeometrie}).

\subsection{Anordnungsstatistik}

Für die Synthese ist nicht nur wichtig, wie die Objekte \emph{aussehen}, sondern
auch, wie sie \emph{liegen}. Tabelle~\ref{tab:anordnung} enthält die dafür
erhobenen Kenngrößen über alle $n = 174$ Labels beider Splits. Alle Abstände sind
in Objektdurchmessern (\Oe) angegeben, sind also skaleninvariant und damit
direkt als Generatorparameter verwendbar.

\begin{table}[htbp]
\centering
\caption[Gemessene Anordnungsstatistik]{Gemessene Anordnungsstatistik der
realen Labels ($n = 174$). NN = nächster Nachbar.}
\label{tab:anordnung}
\footnotesize
\begin{tabular}{@{}lll@{}}
\toprule
\textbf{Größe} & \textbf{\code{val}} & \textbf{\code{test}} \\
\midrule
NN-Abstand (in \Oe) & p10 0{,}86, Median 1{,}09, p90 1{,}78 & p10 0{,}83, Median 1{,}24, p90 2{,}51 \\
Nachbarn innerhalb 5\,\Oe & Median 10, max 14 & Median 10, max 20 \\
berührt NN ($<1{,}2$\,\Oe) & 61{,}4\,\% & 46{,}5\,\% \\
Größenverhältnis zum NN & Median 1{,}05, p90 1{,}14 & Median 1{,}05, p90 1{,}17 \\
Nachbarrichtung & quasi isotrop, leichte Häufung $90$--$150^\circ$ & quasi isotrop \\
\bottomrule
\end{tabular}
\end{table}

\absatz{Befund 8 -- Nachbarn haben nahezu dieselbe Größe.} Das Größenverhältnis
benachbarter Objekte liegt im Median bei 1,05, im 90.\ Perzentil bei 1,14--1,17.
Der Grund ist physikalisch: Innerhalb eines Stapels haben alle Rohre dieselbe
Nennweite; Größenunterschiede entstehen nur durch Perspektive. Der Generator
zieht deshalb \emph{eine Basisgröße pro Cluster} und variiert je Instanz nur um
wenige Prozent. Jede Instanz unabhängig zu würfeln wäre ein sofort sichtbarer
Realismusfehler -- und in der Sprache von
Abschnitt~\ref{sec:shortcut} ein leicht erlernbares Unterscheidungsmerkmal
zwischen synthetischen und echten Bildern.

\absatz{Befund 9 -- jede Stirnfläche hat einen Schaft.} In den Nahaufnahmen gilt
ausnahmslos: Hinter jeder annotierten Stirnfläche führt ein sichtbarer Rohrschaft
weg. Ein dunkler Kreis mitten auf einem durchgehenden Rohr ist physikalisch
unmöglich. Dies ist die stärkste Randbedingung des gesamten Entwurfs und
begründet die aufwendigste Einzelentscheidung der Pipeline
(Abschnitt~\ref{sec:schaft}).

\absatz{Befund 10 -- die Datenlage der Evaluation ist die größte methodische
Schwäche.} Je \emph{ein} Bild für Validierung und Test ist für eine belastbare
Metrikangabe zu wenig: Die Streuung einer mAP-Schätzung auf 88 Objekten ist groß,
und beide Bilder zeigen dieselbe Objektart bei ähnlichem Wetter. Alle berichteten
Zahlen sind entsprechend als Punktschätzungen mit erheblicher Unsicherheit zu
lesen. Das Bewertungsframework gibt deshalb zusätzlich
Bootstrap-Konfidenzintervalle aus (Abschnitt~\ref{sec:bootstrap}).

\section{Was die Messungen festlegen}
\label{sec:messungmapping}

Tabelle~\ref{tab:mapping} fasst zusammen, welche Messung welchen
Pipelineparameter bestimmt. Sie ist die Belegbrücke zwischen diesem Kapitel und
Kapitel~\ref{kap:pipeline} und macht zugleich sichtbar, welche Parameter
\emph{nicht} durch eine Messung gedeckt sind.

\begin{table}[htbp]
\centering
\caption[Messung und abgeleiteter Parameter]{Zuordnung von Messung zu
Pipelineparameter. Die letzte Spalte nennt die Art des Belegs: \emph{M} =
Messung an den Realdaten, \emph{P} = physikalische Randbedingung, \emph{L} =
Literatur, \emph{A} = Annahme.}
\label{tab:mapping}
\footnotesize
\begin{tabularx}{\textwidth}{@{}L{0.30\textwidth}XL{0.055\textwidth}@{}}
\toprule
\textbf{Messung / Befund} & \textbf{Festgelegter Parameter} & \textbf{Art} \\
\midrule
$\sqrt{\text{Fläche}}$ 12,2--35,8\px{} & Objektgrößenbereich des Generators
12--44\px{} (reale Spanne plus 25\,\% Reserve) & M \\
Seitenverhältnis p10--p90 0,85--1,31 & Ovalisierung 57\,\% der Instanzen,
oblat:prolat = 45:12 & M \\
Rohr-\Oe{} der Clean Plates 16--21\px{} & Kopplung Zoom $\leftrightarrow$
Objektgröße & P \\
NN-Abstand Median 1,09--1,24\,\Oe & Gitterabstände des Layoutmodells & M \\
Größenverhältnis zum NN 1,05 & eine Basisgröße pro Cluster, $\pm 8\,\%$ je
Instanz & M \\
5,8\,\% angeschnittene Labels in \code{test} & Clusterzentrum darf außerhalb der
Leinwand liegen; Randregel 15\,\% & M \\
Befund 9 (Schaft hinter jedem Profil) & Rohrschaft bei 88\,\% der Instanzen & P \\
Befund 5 (Maschendraht, Schotter) & 10\,\% Negativbilder & M/L \\
Befund 1 (Alphakanal unbrauchbar) & eigene Matting-Stufe & P \\
Befund 2 (dunkle/helle Bohrung) & Gleichverteilung über alle neun Ausschnitte & A \\
keine Messung verfügbar & Anteil verformter Profile (\enquote{Dellen}) 18\,\% & A \\
keine Messung verfügbar & Anteil Schrägansicht (Perspektive) 30\,\% & A \\
\bottomrule
\end{tabularx}
\end{table}

Zwei Zeilen der Tabelle sind ausdrücklich als Annahme gekennzeichnet. In den
Realbildern sind Dellen und Schrägansichten nicht separat annotiert, ihre
Häufigkeit ist also unbekannt. Nach Anforderung~A6 werden die gewählten Werte
dennoch benannt und ihre Wirkungsrichtung diskutiert
(Abschnitte~\ref{sec:instanzgeometrie} und \ref{sec:grenzen}), statt sie als
gemessen auszugeben.
